%!TEX root = ../main.tex
\chapter*{Glossario}
\addcontentsline{toc}{chapter}{Glossario}

I termini tecnici ricorrenti nella tesi sono definiti qui una volta sola, in
ordine alfabetico.

\begin{description}
  \item[Adapter] Componente che avvolge uno strumento esterno, ne esegue
  l'analisi e ne traduce l'uscita nel modello di dominio del sistema, così che
  il resto dell'applicazione non dipenda dal formato del singolo strumento.

  \item[Albero sintattico] (\emph{syntax tree}) Rappresentazione strutturata del
  codice sorgente prodotta da un parser, in cui funzioni, istruzioni ed
  espressioni sono nodi di un albero; analizzarlo, anziché il testo, rende
  l'esame indifferente a commenti e formattazione.

  \item[Allowlist e denylist] (anche \emph{whitelist} e \emph{blacklist}) Elenco
  dei valori ammessi e, rispettivamente, dei valori esclusi. Filtrare con
  un'allowlist blocca tutto ciò che non è previsto; una denylist lascia passare
  qualunque valore non elencato, compresi quelli aggiunti in seguito.

  \item[Analisi della contaminazione] (\emph{taint analysis}) Tecnica che marca
  come ``contaminati'' i valori provenienti da una \emph{sorgente} esterna, cioè
  il punto in cui un dato controllato dall'esterno entra nel programma, e ne
  segue la propagazione fino a un \emph{sink}, l'operazione sensibile in cui il
  valore viene infine usato.

  \item[Analisi statica e dinamica] L'analisi statica esamina gli artefatti di
  progetto (codice, configurazioni, contratti) senza eseguire il sistema;
  l'analisi dinamica invia richieste al sistema in esecuzione e ne osserva le
  risposte.

  \item[API ombra] (\emph{shadow API}) Endpoint esposto dall'applicazione ma
  assente dalla documentazione ufficiale, e quindi sottratto ai controlli e
  alla manutenzione di chi la gestisce.

  \item[Autenticazione e autorizzazione] L'autenticazione accerta chi effettua
  la richiesta; l'autorizzazione stabilisce se quel soggetto, una volta
  identificato, abbia il diritto di compiere una specifica operazione su una
  specifica risorsa.

  \item[Black box] Modalità di analisi che osserva il sistema solo dall'esterno,
  attraverso le risposte alle richieste, senza accesso al codice sorgente; vi si
  contrappone l'analisi \emph{white box}, che legge il codice.

  \item[Bucket] Contenitore di oggetti in un servizio di storage cloud come
  Amazon S3, con permessi di accesso propri che ne determinano la visibilità
  pubblica o privata.

  \item[\texttt{Cache-Control}] Intestazione HTTP che indica a browser e proxy se
  e per quanto tempo possono conservare una copia della risposta; se manca, dati
  riservati possono restare in una cache condivisa.

  \item[Chiave API] (\emph{API key}) Stringa segreta inviata con ogni richiesta
  per identificare l'applicazione chiamante; identifica un client, non un
  utente, e non sostituisce un meccanismo di autenticazione.

  \item[Claim] Singola affermazione contenuta in un token, per esempio
  l'identificativo dell'utente (\texttt{sub}), il suo ruolo, la scadenza o il
  destinatario previsto (\emph{audience}).

  \item[Clean Architecture] Stile architetturale che organizza il software in
  livelli concentrici (dominio, applicazione, infrastruttura, presentazione), in
  cui le dipendenze puntano sempre verso il dominio e i dettagli tecnici restano
  ai margini.

  \item[Codice di stato HTTP] Numero di tre cifre con cui il server riassume
  l'esito di una richiesta, per esempio \texttt{200 OK} (successo),
  \texttt{403 Forbidden} (accesso negato), \texttt{404 Not Found} (risorsa
  inesistente), \texttt{429 Too Many Requests} (limite di frequenza superato).

  \item[Confine di fiducia] (\emph{trust boundary}) La linea che separa la parte
  di un sistema di cui ci si fida da quella che va invece verificata;
  collocarlo nel punto sbagliato significa accettare senza controlli dati che
  provengono in realtà dall'esterno.

  \item[Confusione di algoritmo] (\emph{algorithm confusion}) Attacco ai token
  firmati in cui l'attaccante cambia l'algoritmo dichiarato, per esempio da
  RS256, asimmetrico, a HS256, simmetrico, inducendo il server a verificare la
  firma con una chiave che l'attaccante conosce.

  \item[Container] Unità di esecuzione isolata che impacchetta un'applicazione
  con le sue dipendenze; condivide il kernel del sistema ospite, per cui un
  container eseguito con privilegi elevati ne amplia la superficie di attacco.

  \item[Controller] Componente di un framework web che raggruppa le funzioni
  incaricate di gestire le richieste rivolte a una stessa risorsa.

  \item[CORS] (\emph{Cross-Origin Resource Sharing}) Meccanismo con cui un server
  dichiara quali siti di origine diversa possono leggere le sue risposte dal
  browser; una configurazione permissiva consente a domini terzi non
  autorizzati di accedere a dati riservati.

  \item[CVSS] (\emph{Common Vulnerability Scoring System}) Standard pubblico per
  attribuire a una vulnerabilità un punteggio di gravità da 0 a 10, calcolato a
  partire da un vettore di metriche sulle sue caratteristiche.

  \item[CWE] (\emph{Common Weakness Enumeration}) Il catalogo pubblico delle
  debolezze software mantenuto da MITRE; ogni voce, indicata dalla sigla CWE
  seguita da un numero, descrive una classe di difetto indipendentemente dal
  prodotto in cui compare.

  \item[Decoratore] Costrutto di linguaggi come Python che avvolge una funzione
  per aggiungerle un comportamento; nei framework web è il modo tipico di
  proteggere una rotta, per esempio con \texttt{login\_required}.

  \item[Denial of service] (DoS) Attacco che rende un servizio indisponibile
  agli utenti legittimi, per esempio esaurendone CPU, memoria o banda.

  \item[Deny-by-default] (diniego per impostazione predefinita) Il criterio per
  cui ogni funzione è negata salvo concessione esplicita al ruolo che ne ha
  diritto; con il criterio opposto una dimenticanza non produce alcun errore e
  passa inosservata.

  \item[Differential testing] Tecnica di verifica che invia la stessa richiesta
  con identità diverse e ne confronta le risposte: se un utente ottiene ciò che
  spetta a un altro, il controllo di autorizzazione è assente o difettoso.

  \item[Directory listing] Funzione di un server web che, in assenza di una
  pagina indice, mostra l'elenco dei file contenuti in una cartella, rivelando
  risorse che non dovrebbero essere pubbliche.

  \item[DTO] (\emph{Data Transfer Object}, oggetto di trasporto) Oggetto usato
  per trasferire dati fra il client e il server, spesso una proiezione parziale
  dell'entità memorizzata nel database.

  \item[Endpoint] Il singolo punto di accesso di un'API, identificato da un URL e
  da un metodo HTTP, che espone una specifica operazione su una risorsa.

  \item[Entropia di Shannon] Misura del grado di casualità di una stringa;
  sequenze ad alto disordine indicano chiavi o token reali, distinguendoli dai
  normali identificatori.

  \item[Escalation dei privilegi] Il passaggio di un utente a permessi superiori
  a quelli che gli sono stati concessi, per esempio da utente ordinario ad
  amministratore.

  \item[Falso positivo e falso negativo] Un falso positivo è una segnalazione
  che non corrisponde a un difetto reale; un falso negativo è un difetto reale
  che lo strumento non segnala. Vi si contrappone il \emph{vero positivo}, una
  segnalazione che corrisponde a un difetto reale.

  \item[Finding] La singola segnalazione prodotta da uno strumento di analisi,
  tradotta nel modello di dominio comune del sistema con sorgente, posizione,
  categoria, severità e confidenza.

  \item[Fingerprinting del dispositivo] Tecnica che riconosce un client a partire
  da caratteristiche del dispositivo e del browser, anche quando cambia
  indirizzo IP o account.

  \item[Forza bruta] (\emph{brute force}) Attacco che prova sistematicamente un
  gran numero di valori, per esempio password o codici di verifica, finché non
  trova quello corretto.

  \item[Fuzzing] Tecnica di verifica che invia a un sistema un gran numero di
  input generati automaticamente, spesso malformati o inattesi, osservandone
  le reazioni anomale.

  \item[Gateway] Componente posto davanti ai servizi applicativi che riceve
  tutte le chiamate esterne e vi applica funzioni comuni quali autenticazione,
  instradamento e limiti di frequenza.

  \item[GraphQL] Linguaggio di interrogazione per API in cui il client, anziché
  richiamare endpoint predefiniti, descrive in una sola richiesta l'esatto
  insieme di dati che desidera ricevere.

  \item[Handler] La funzione che il framework invoca per servire una determinata
  rotta: riceve la richiesta, esegue l'operazione e costruisce la risposta.

  \item[Hardening] (irrobustimento) Insieme delle operazioni con cui si riduce la
  superficie di attacco di un sistema: disattivazione di funzioni e servizi non
  necessari, rimozione delle credenziali predefinite, applicazione delle
  impostazioni consigliate.

  \item[HSTS] (\emph{HTTP Strict Transport Security}) Intestazione HTTP con cui un
  server impone al browser di contattarlo soltanto su connessione cifrata.

  \item[IaC] (\emph{Infrastructure as Code}) Pratica per cui le risorse
  infrastrutturali (reti, istanze, storage, permessi) sono dichiarate in file
  versionati che uno strumento di provisioning traduce nello stato reale; il
  file diventa così un artefatto ispezionabile prima che la risorsa esista.

  \item[Identity provider] Servizio esterno che gestisce le identità degli utenti,
  ne verifica le credenziali e rilascia i token con cui le applicazioni li
  riconoscono, per esempio Keycloak.

  \item[In chiaro] (\emph{cleartext}) Senza cifratura; un URL \texttt{http://}
  anziché \texttt{https://} espone i dati in transito a chi intercetta il
  traffico, l'attacco \emph{man-in-the-middle}.

  \item[Iniezione] (\emph{injection}) Attacco in cui un dato fornito
  dall'esterno viene interpretato come comando, per esempio come parte di una
  query SQL (\emph{SQL injection}). Nell'iniezione \emph{differita} il dato
  malevolo viene prima memorizzato e solo in un secondo momento usato.

  \item[Integrazione continua] (\emph{continuous integration}) Pratica per cui a
  ogni modifica del codice una pipeline automatica compila, verifica e analizza
  il progetto.

  \item[Intestazione HTTP] (\emph{header}) Campo di metadati che accompagna una
  richiesta o una risposta HTTP, per esempio \texttt{Authorization}, che
  trasporta le credenziali, o \texttt{Content-Type}, che indica il formato del
  corpo.

  \item[JWKS] (\emph{JSON Web Key Set}) Insieme delle chiavi pubbliche, pubblicato
  in formato JSON, con cui un servizio verifica la firma dei token; se il server
  accetta un riferimento a un JWKS indicato nel token stesso, l'attaccante può
  firmare token con una chiave propria.

  \item[JWT] (\emph{JSON Web Token}) Formato di token composto da tre parti
  codificate in Base64 e separate da punti: intestazione, contenuto (i
  \emph{claim}) e firma. La firma garantisce l'integrità del contenuto solo se
  il server la verifica.

  \item[Least privilege] (principio del minimo privilegio) Il principio per cui
  un soggetto dispone soltanto dei permessi strettamente necessari a svolgere
  le proprie funzioni.

  \item[Lista di controllo d'accesso] (ACL, \emph{Access Control List}) Elenco
  associato a una risorsa che indica quali soggetti possono accedervi e con
  quali operazioni.

  \item[Log4Shell] Vulnerabilità del 2021 della libreria di logging Java Log4j,
  in cui una stringa registrata nei log poteva indurre il server a scaricare ed
  eseguire codice remoto.

  \item[Manifest] File dichiarativo che descrive una risorsa o un progetto: un
  manifest Terraform dichiara risorse infrastrutturali, un manifest delle
  dipendenze (come \path{requirements.txt} o \path{package.json}) elenca le
  librerie usate.

  \item[Mass assignment] L'assegnazione automatica dei campi ricevuti dal client
  agli attributi di un oggetto applicativo; se il server non filtra i campi
  ammessi, il client può valorizzare attributi interni non previsti dal flusso
  normale.

  \item[Metodo HTTP] (o \emph{verbo HTTP}) L'azione richiesta su una risorsa, per
  esempio \texttt{GET} per leggere o \texttt{DELETE} per cancellare.

  \item[MFA] (\emph{Multi-Factor Authentication}) Autenticazione che richiede,
  oltre alla password, un secondo fattore indipendente, spesso un codice
  monouso (OTP, \emph{One-Time Password}).

  \item[Middleware] Componente che si interpone fra la ricezione di una
  richiesta e il suo handler, applicando funzioni comuni a più rotte come
  l'autenticazione o la registrazione degli accessi.

  \item[Modello linguistico] (\emph{Large Language Model}) Modello di
  intelligenza artificiale addestrato su grandi quantità di testo, capace di
  generare risposte in linguaggio naturale; nel sistema è eseguito localmente
  tramite Ollama.

  \item[Multi-tenant] Architettura in cui più clienti (\emph{tenant}) condividono
  la stessa infrastruttura applicativa, separati soltanto dalla logica di
  autorizzazione.

  \item[OAuth e OpenID Connect] OAuth~2.0 è il protocollo con cui un'applicazione
  ottiene un accesso delegato alle risorse di un utente; OpenID Connect (OIDC)
  vi aggiunge l'autenticazione dell'utente. Parametri come \texttt{state} e PKCE
  (\emph{Proof Key for Code Exchange}) proteggono il flusso da intercettazioni e
  falsificazioni.

  \item[OpenAPI] Formato standard, in YAML o JSON, per descrivere un'API REST:
  endpoint, parametri, schemi di richiesta e risposta, requisiti di sicurezza.
  Il documento che ne risulta è detto \emph{contratto} o \emph{specifica}
  dell'API.

  \item[Operazione batch] Una singola richiesta che trasporta un elenco di
  elementi da elaborare insieme.

  \item[Paginazione] La suddivisione di una risposta voluminosa in pagine di
  dimensione limitata, così che una singola richiesta non restituisca
  un'intera collezione.

  \item[Parametro di percorso e di query] Valori trasportati dall'URL di una
  richiesta: il parametro di percorso ne fa parte (\texttt{482} in
  \texttt{/orders/482}), quello di query segue il punto interrogativo
  (\texttt{?limit=50}).

  \item[Payload] Il contenuto utile trasportato da una richiesta o da una
  risposta, tipicamente il corpo in formato JSON; per estensione, il contenuto
  costruito da un attaccante per sfruttare una vulnerabilità.

  \item[Pipeline] Sequenza automatizzata di fasi eseguite una dopo l'altra, in
  cui l'uscita di ciascuna alimenta la successiva.

  \item[Policy] Documento che dichiara quali soggetti (\emph{principal}) possono
  eseguire quali azioni su quali risorse; nei servizi cloud governa, per
  esempio, l'accesso a un bucket o i permessi di un ruolo IAM.

  \item[Prova d'unità] (\emph{unit test}) Verifica automatica di un singolo
  componente del software isolato dal resto, con le dipendenze esterne
  sostituite da simulazioni.

  \item[Proxy] Componente intermedio attraverso cui transitano le richieste fra
  client e server; un proxy di intercettazione permette di osservare e
  modificare il traffico.

  \item[Rate limiting] (limitazione di frequenza) Il limite al numero di
  richieste che un client può inviare in un dato intervallo di tempo.

  \item[Recall e precisione] Metriche di accuratezza di uno strumento di
  rilevamento: la recall è la quota di difetti reali che lo strumento trova, la
  precisione la quota di segnalazioni che corrispondono a difetti reali.

  \item[Redirezione] (\emph{redirect}) La risposta con cui un server rimanda il
  client a un altro indirizzo; seguirla in automatico consente a un fornitore
  compromesso di dirottare la richiesta, dati sensibili compresi, verso un
  server ostile.

  \item[ReDoS] (\emph{Regular expression Denial of Service}) Un input costruito
  ad arte che fa degenerare il tempo di valutazione di un'espressione regolare,
  saturando la CPU con una richiesta sola.

  \item[Refresh token] Token a lunga durata usato per ottenere nuovi token di
  accesso senza ripetere il login; per limitarne l'abuso dovrebbe essere
  monouso, ruotato a ogni utilizzo e soggetto a scadenza.

  \item[Remediation] L'insieme delle azioni correttive che eliminano o mitigano
  una vulnerabilità segnalata.

  \item[Riferimento di verità] (\emph{ground truth}) L'elenco delle
  vulnerabilità effettivamente presenti in un bersaglio, rispetto al quale si
  misurano le segnalazioni di uno strumento.

  \item[Rotta] (\emph{route}) L'associazione fra un percorso URL con il relativo
  metodo HTTP e la funzione che lo gestisce; il \emph{routing} è il meccanismo
  che smista ogni richiesta verso la rotta corrispondente.

  \item[Runtime] La fase in cui il software è in esecuzione, contrapposta a
  quella in cui viene scritto, compilato o analizzato staticamente.

  \item[SAML e SSO] Il \emph{Single Sign-On} (SSO) consente di accedere a più
  servizi con un'unica autenticazione; SAML (\emph{Security Assertion Markup
  Language}) è uno standard XML per realizzarlo, in cui l'identity provider
  rilascia asserzioni firmate sull'identità dell'utente.

  \item[Schema di sicurezza] (\emph{security scheme}) La sezione del contratto
  OpenAPI che dichiara come un endpoint è protetto, per esempio tramite token o
  chiave API.

  \item[Serializzazione] La conversione di un oggetto in memoria in un formato
  trasmissibile, come JSON; serializzare automaticamente un oggetto intero può
  esporne anche gli attributi interni.

  \item[Serverless] Modello di esecuzione cloud in cui il codice gira in funzioni
  attivate su richiesta, senza che lo sviluppatore gestisca i server
  sottostanti.

  \item[Servizio di metadati] (\emph{instance metadata service}) Servizio
  interno degli ambienti cloud, raggiungibile da ogni macchina a un indirizzo
  noto come \texttt{169.254.169.254}, che fornisce configurazione e spesso
  credenziali temporanee dell'istanza.

  \item[Session fixation] Attacco in cui l'avversario impone alla vittima un
  identificativo di sessione da lui conosciuto; se il server non lo rinnova
  dopo il login, l'attaccante eredita la sessione autenticata.

  \item[Spoofing] La falsificazione di un'informazione per farsi passare per
  qualcun altro, per esempio un'intestazione HTTP che dichiara un indirizzo IP
  di origine diverso da quello reale.

  \item[Terraform] Strumento di Infrastructure as Code che legge file
  dichiarativi in linguaggio HCL e crea, modifica o distrugge le risorse cloud
  corrispondenti.

  \item[Timeout] Il tempo massimo concesso a un'operazione prima di
  interromperla; in sua assenza una chiamata lenta può restare appesa a tempo
  indefinito, trattenendo risorse.

  \item[TLS] (\emph{Transport Layer Security}) Il protocollo che cifra e
  autentica il canale di comunicazione su cui viaggia il traffico HTTP.

  \item[Token] Stringa rilasciata dopo l'autenticazione e inviata con ogni
  richiesta successiva, tipicamente nell'intestazione
  \texttt{Authorization: Bearer}, per dimostrare l'identità del chiamante senza
  ripetere le credenziali.

  \item[Traccia dello stack] (\emph{stack trace}) L'elenco delle chiamate di
  funzione attive nel momento dell'errore; mostrata al client rivela percorsi,
  librerie e struttura interna dell'applicazione.

  \item[UUID] (\emph{Universally Unique Identifier}) Identificatore di 128 bit
  generato in modo da essere praticamente unico e, a differenza di un intero
  progressivo, difficile da indovinare.

  \item[Versioning] La conservazione delle versioni precedenti di un oggetto,
  che consente di recuperarlo dopo una modifica o una cancellazione.

  \item[Webhook] Indirizzo che un'applicazione registra presso un servizio
  perché questo vi invii una richiesta HTTP al verificarsi di un evento.
\end{description}
\clearpage
